\BourbakiExerciseGroup

\begin{enumerate}
\def\labelenumi{\arabic{enumi})}
\item
  Let X, Y be two topological spaces, A a subset of X and B a subset of Y. Show that \(\mathrm{Fr}(\mathbf{A} \times \mathbf{B}) = (\mathrm{Fr}(\mathbf{A}) \times \overline{\mathbf{B}}) \cup (\overline{\mathbf{A}} \times \mathrm{Fr}(\mathbf{B}))\) .
\item
  Let \(\mathbf{X}\) , \(\mathbf{Y}\) be two topological spaces, A a closed subset of \(\mathbf{X} \times \mathbf{Y}\) . If \(\mathfrak{B}(x)\) is the set of all neighbourhoods of a point \(x \in \mathbf{X}\) , show that
\end{enumerate}

\[
\bigcap_ {\mathbf {V} \in \mathfrak {B} (x)} \overline {{\mathrm{A} (\mathbf {V})}} = \mathrm{A} (x).
\]

Give an example of an open subset B of a product \(X \times Y\) such that

\[
\bigcap_ {\mathbf {V} \in \mathfrak {B} (x)} \overline {{\mathbf {B} (\mathbf {V})}}
\]

is different from \(\mathbf{B}_{(x)}\) .

\begin{enumerate}
\def\labelenumi{\arabic{enumi})}
\setcounter{enumi}{2}
\tightlist
\item
  Let \((\mathbf{X}_t)_{i\in I}\) be an infinite family of topological spaces such that each \(\mathbf{X}_t\) contains (at least) two distinct points \(a_i, b_i\) such that there is a neighbourhood of \(b_i\) which does not contain \(a_i\) .
\end{enumerate}

\begin{enumerate}
\def\labelenumi{\alph{enumi})}
\tightlist
\item
  For each index \(x \in I\) , let \(c_{x}\) be the point of the product space
\end{enumerate}

\[
\mathbf {X} = \prod_ {k + 1} \mathbf {X} _ {i}
\]

for which \(\mathrm{pr}_x c_x = b_x\) and \(\mathrm{pr}_t c_x = a_t\) whenever \(i \neq x\) . Show that every point of the set C of points \(c_x\) is isolated.

\begin{enumerate}
\def\labelenumi{\alph{enumi})}
\setcounter{enumi}{1}
\item
  Deduce from a) that the topology of X has a countable base if and only if I is countable and the topology of each \(X_{i}\) has a countable base (use Exercise 7 of §1).
\item
  Show that if the index set I is not countable then the point \(b = (b_{i})\) of X has no countable fundamental system of neighbourhoods.
\end{enumerate}

\begin{enumerate}
\def\labelenumi{\arabic{enumi})}
\setcounter{enumi}{3}
\tightlist
\item
  Let K be the discrete space consisting of the two numbers 0,1; let A be an infinite set, and let X be the product space \(\mathbf{K}^{\mathbf{A}}\) . Let
\end{enumerate}

\[
\mathbf {V} = \prod_ {\alpha \in \mathbf {A}} \mathbf {V} _ {\alpha}
\]

be an elementary set in X; if h is the number of indices \(\alpha\) such that \(V_{\alpha} \neq K\) , put \(\mu(V) = 2^{-h}\) .

\begin{enumerate}
\def\labelenumi{\alph{enumi})}
\tightlist
\item
  Show that if \(\mathbf{U}_1, \ldots, \mathbf{U}_n\) are pairwise disjoint elementary sets, then
\end{enumerate}

\[
\sum_ {k = 1} ^ {n} \mu (\mathrm{U} _ {k}) \leqslant \mathrm{I}.
\]

{[}Write each \(U_{k}\) in the form \(W_{k} \times K^{B}\) , where B (independent of k) is the complement of a finite subset of A, and \(W_{k}\) is a finite set; count up the number of points in \(W_{k}\) {]}.

\begin{enumerate}
\def\labelenumi{\alph{enumi})}
\setcounter{enumi}{1}
\item
  Deduce from a) that X satisfies the condition \((\mathbf{D}_{\mathbf{IV}})\) of Exercise 7 of §1. {[}If there were an uncountable family \((\mathbf{U}_{\lambda})\) of pairwise disjoint elementary sets, there would exist an integer p \textgreater{} 0 such that \(\mu(\mathbf{U}_{\lambda}) \geqslant 1/p\) for an infinite number of indices \(\lambda\) {]}.
\item
  Show that if A is uncountable, then X does not satisfy the condition \((\mathbf{D}_{\mathrm{III}})\) of Exercise 7 of §1 (cf.~Exercise 3).
\end{enumerate}

\begin{enumerate}
\def\labelenumi{\arabic{enumi})}
\setcounter{enumi}{4}
\tightlist
\item
  Let K be the discrete space consisting of the two numbers 0, 1; let A be an infinite set and let X denote the product space \(\mathbf{K}^{\mathfrak{P}(\mathbf{A})}\) . a) Let J be any finite subset of A; suppose J has p elements. For each subset H of J, let \(M_{H}\) be the set of subsets L of A such that \(L \cap J = H\) ; the \(M_{H}\) form a partition \(\tilde{\omega}_{J}\) of \(\mathfrak{P}(A)\) into \(2^{p}\) subsets. Let \(F_{J}\) be the subset of X consisting of those points \((x_{L})_{L \subset A}\) such that \(x_{L} = x_{M}\) whenever L and M belong to the same set of the partition \(\tilde{\omega}_{J}\) . \(F_{J}\) is a finite set, with \(2^{2^{p}}\) elements. If F is the union of the \(F_{J}\) as J runs through all finite subsets of A, show that F is equipotent to A.
\end{enumerate}

\begin{enumerate}
\def\labelenumi{\alph{enumi})}
\setcounter{enumi}{1}
\tightlist
\item
  Show that F is dense in X. Taking A = N, deduce that X satisfies condition \((D_{\mathrm{II}})\) of §1, Exercise 7, but does not satisfy condition \((D_{\mathrm{III}})\) (Exercise 4).
\end{enumerate}

\begin{enumerate}
\def\labelenumi{\arabic{enumi})}
\setcounter{enumi}{5}
\item
  Let X, Y, Z be three topological spaces, A an open subset of \(X \times Y\) and B an open subset of \(Y \times Z\) . Show that \(B \circ A\) is an open subset of \(X \times Z\) .
\item
  Let X be the sum of a family \((X_{\lambda})\) of topological spaces, Y the sum of a family \((Y_{\mu})\) of topological spaces. Show that the product space \(X \times Y\) is the sum of the \(X_{\lambda} \times Y_{\mu}\) .
\item
  Let \((\mathbf{G}_t)_{t \in \mathbf{I}}\) be a family of graphs of equivalence relations (Set Theory, Chapter II, § 6, no. 1) on a topological space \(\mathbf{X}\) , and let \(\mathbf{R}_t\) be the relation \((x, y) \in \mathbf{G}_t\) ; then
\end{enumerate}

\[
\mathbf {G} = \bigcap_ {\iota \in \mathbf {I}} \mathbf {G} _ {\iota}
\]

is the graph of an equivalence relation R on X. Show that there is a canonical continuous injection of the quotient space X/R into the product space \(\prod_{i}\left(X/R_{i}\right)\) .

\begin{enumerate}
\def\labelenumi{\arabic{enumi})}
\setcounter{enumi}{8}
\tightlist
\item
  Let \((\mathbf{X}_i)_{i\in I}\) be an infinite family of topological spaces. Study the topology on the product set
\end{enumerate}

\[
\mathbf {X} = \prod_ {i \in I} \mathbf {X} _ {i} \text {   generated   by   subsets   of   the   form   } \prod_ {i \in I} \mathbf {U} _ {i},
\]

where \(U_{t}\) is an open subset of \(X_{t}\) for each \(t \in I\) and the set of indices t for which \(U_{t} \neq X_{t}\) has cardinal less than c, where c is a given infinite cardinal. Consider what becomes of the propositions of §4 for this topology.

\begin{enumerate}
\def\labelenumi{\arabic{enumi})}
\setcounter{enumi}{9}
\tightlist
\item
  \begin{enumerate}
  \def\labelenumii{\alph{enumii})}
  \tightlist
  \item
    Let \((\mathbf{X}_{\alpha}, f_{\alpha\beta})\) be an inverse system of topological spaces, and let X be the inverse limit. Let \(f_{\alpha}\) be the canonical mapping \(X \to X_{\alpha}\) . For each \(\alpha\) , let \(Y_{\alpha}\) be a subspace of \(X_{\alpha}\) containing \(f_{\alpha}(X)\) , such that the \(Y_{\alpha}\) form an inverse system of subspaces of the \(X_{\alpha}\) . Show that \(\varprojlim Y_{\alpha}\) can be canonically identified with X.
  \end{enumerate}
\end{enumerate}

\begin{enumerate}
\def\labelenumi{\alph{enumi})}
\setcounter{enumi}{1}
\tightlist
\item
  Let \((\mathbf{X}_{\alpha}^{\prime}, f_{\alpha\beta}^{\prime})\) be another inverse system of topological spaces with the same directed index set. For each \(\alpha\) , let \(u_{\alpha}: X_{\alpha} \to X_{\alpha}^{\prime}\) be a continuous mapping such that the \(u_{\alpha}\) form an inverse system of mappings. The \(u_{\alpha}(\mathbf{X}_{\alpha})\) form an inverse system of subspaces of the \(X_{\alpha}^{\prime}\) ; show that if \(u = \varprojlim u_{\alpha}\) and if the \(f_{\alpha}\) are surjective, then \(u(\mathbf{X})\) is dense in the space \(\varprojlim u_{\alpha}(\mathbf{X}_{\alpha})\) . Does this result still hold good if the \(f_{\alpha}\) are no longer assumed to be surjective? (Cf. Set Theory, Chapter III, §1, Exercise 32.)
\end{enumerate}

